\documentclass[11pt]{ctexart}
\usepackage[a4paper,top=2cm,bottom=2cm,left=2.2cm,right=2.2cm]{geometry}
\usepackage{fancyvrb}
\usepackage{xcolor}
\usepackage{colortbl}
\usepackage{fontspec}
\usepackage[most]{tcolorbox}
\usepackage{booktabs}
\usepackage{enumitem}
\usepackage{titlesec}
\usepackage{hyperref}

\setmonofont{Consolas}
\setCJKmonofont{SimHei}
\setlength{\parindent}{0pt}
\setlist{nosep,leftmargin=1.4em}

% ── 上板系列主题色：深松绿（真机/硬件感） ──
\definecolor{deep}{RGB}{32,72,40}        % 主色：深松绿
\definecolor{accent}{RGB}{74,122,66}     % 强调：松绿
\definecolor{keyline}{RGB}{90,130,200}   % 知识点：蓝
\definecolor{keybg}{RGB}{238,244,255}
\definecolor{warnline}{RGB}{200,120,60}  % 警告：橙
\definecolor{warnbg}{RGB}{255,243,235}
\definecolor{okline}{RGB}{26,116,88}     % 成果：绿
\definecolor{okbg}{RGB}{240,252,244}
\definecolor{tipline}{RGB}{110,150,110}  % 提示：浅松绿
\definecolor{tipbg}{RGB}{240,248,240}
\definecolor{codebg}{RGB}{249,249,252}   % 代码底
\definecolor{rowtint}{RGB}{238,246,238}  % 表头浅绿

\titleformat{\section}{\large\bfseries\color{deep}}{\thesection}{0.6em}{}
\titleformat{\subsection}{\normalsize\bfseries\color{accent}}{\thesubsection}{0.5em}{}
\titlespacing{\section}{0pt}{14pt}{6pt}
\titlespacing{\subsection}{0pt}{10pt}{4pt}

\newtcolorbox{keybox}[1]{enhanced, colback=keybg, colframe=keyline,
  boxrule=0.6pt, arc=1.5mm, left=3mm, right=3mm, top=2mm, bottom=2mm,
  title={\bfseries #1}, fonttitle=\small\bfseries, coltitle=white,
  colbacktitle=keyline}
\newtcolorbox{warnbox}[1]{enhanced, colback=warnbg, colframe=warnline,
  boxrule=0.6pt, arc=1.5mm, left=3mm, right=3mm, top=2mm, bottom=2mm,
  title={\bfseries #1}, fonttitle=\small\bfseries, coltitle=white,
  colbacktitle=warnline}
\newtcolorbox{okbox}[1]{enhanced, colback=okbg, colframe=okline,
  boxrule=0.6pt, arc=1.5mm, left=3mm, right=3mm, top=2mm, bottom=2mm,
  title={\bfseries #1}, fonttitle=\small\bfseries, coltitle=white,
  colbacktitle=okline}
\newtcolorbox{tipbox}[1]{enhanced, colback=tipbg, colframe=tipline,
  boxrule=0.5pt, arc=1.5mm, left=3mm, right=3mm, top=1.5mm, bottom=1.5mm,
  title={\bfseries #1}, fonttitle=\small\bfseries, coltitle=white,
  colbacktitle=tipline}

\newcommand{\code}[1]{\texttt{#1}}
\DefineVerbatimEnvironment{Verilog}{Verbatim}{frame=single,fontsize=\small}

\begin{document}

% ==================== 封面 ====================
\begin{titlepage}
\centering
\vspace*{2.2cm}
{\color{deep}\Huge\bfseries 视频通路使用指导}\\[6pt]
{\color{deep}\Large\bfseries 导入 · 仿真 · 烧录（SC520CS）}\\[14pt]
{\large 2026 FPGA 创新设计赛道（安路）赛题二}\\[6pt]
{\large 上板阶段 · 08\_board\_bringup/01\_video\_passthrough}\\[20pt]

\begin{tcolorbox}[enhanced, colback=codebg, colframe=accent, boxrule=0.8pt,
  arc=2mm, width=12.5cm, left=4mm, right=4mm, top=3mm, bottom=3mm]
\centering
{\small 读者：拿到摄像头与开发板的全体队员\\
目标：跑通"SC520CS 摄像头画面到屏幕"的完整通路\\
工程：SC520CS 适配版官方例程 Lab1（1280x720@60fps）\\
配套：ModelSim 整链仿真（虚拟传感器）+ 本指导 PDF}
\end{tcolorbox}

\vspace{8mm}
\begin{tcolorbox}[enhanced, colback=tipbg, colframe=tipline, boxrule=0.5pt,
  arc=2mm, width=12.5cm, left=4mm, right=4mm, top=2mm, bottom=2mm]
\centering
{\small\color{deep} \textbf{使用顺序建议}\\
先读第 1～3 章认识链路与工程 $\rightarrow$ 第 4 章 TD 编译出 bit
$\rightarrow$ 第 6 章烧录上电验收；\\
第 5 章仿真可在等编译时并行做（不需要板子）。}
\end{tcolorbox}

\vfill
{\small\color{black!50} fpga-ptz-tracker 教程系列 · 2026-10}
\end{titlepage}

\tableofcontents
\newpage

% ============================================================
\section{这一步在干什么}

\subsection{链路全景}

\begin{Verilog}
SC520CS 摄像头 (1280x720@60, RAW10, 2-lane MIPI)
   |  FPC 排线
   v
[MIPI D-PHY 接收]  -> [CSI 解包] -> [BLC 黑电平] -> [demosaic 去马赛克]
                                                    |
                       [HDMI 发送 (TMDS)] <- [mixer/VTC] <- [AWB 白平衡]
                              |
                              v
                          显示器 720p60
\end{Verilog}

中间还有 DDR2 帧缓存（写读跨时钟域）。整个工程是官方例程 Lab1 的
SC520CS 适配版（厂商已把 Lab1/2/3 统一配置为 SC520CS 原厂
1280x720@60fps HD 模式），本步\textbf{不加任何视觉算法}——
先把"画面能上屏"这一段在真机上点亮。

\subsection{与前阶段的关系}

\begin{itemize}
  \item 前五个 Step 的视觉算法（检测/OSD/云台）都是在这条通路的
    96bit 节点上"窃流"挂载的——本步通路是它们的宿主；
  \item 嵌入式教程（一）讲的 \code{system\_top} 集成、行为替身，
    就是围绕这条通路做的仿真版；本步把它换成真硬件；
  \item 通路上板成功后，按嵌入式教程（二）的 runbook 固化，
    再逐级挂视觉模块（五步联调节奏）。
\end{itemize}

% ============================================================
\section{硬件连接}

\subsection{连接清单}

\begin{center}
\begin{tabular}{lll}
\toprule
\rowcolor{rowtint}\textbf{连接} & \textbf{接口} & \textbf{注意} \\
\midrule
摄像头模组 & 板上 CAM FPC 座 & 排线\textbf{金手指方向}按座子丝印，勿反插 \\
HDMI 显示器 & 板上 HDMI TX 座 & 显示器需支持 1280x720@60 \\
电源 & 板上电源口 & 先接显示器再上电（或按板卡手册顺序） \\
下载器 & JTAG/USB & 调试下载与固化共用 \\
\bottomrule
\end{tabular}
\end{center}

\begin{warnbox}{反插排线是第一大硬件事故源}
FPC 排线插反可能短路烧摄像头。插线前：断电、对照座子丝印的
触点位置与排线蓝面/金手指方向。插好后再上电。
\end{warnbox}

\subsection{本工程用到的板级引脚（pin.adc 摘录）}

\begin{center}
\begin{tabular}{lll}
\toprule
\rowcolor{rowtint}\textbf{信号} & \textbf{引脚} & \textbf{作用} \\
\midrule
\code{I\_sys\_clk} & C4 & 系统时钟 \\
\code{I\_rst\_n} & R10 & 复位（低有效） \\
\code{O\_cam\_scl} / \code{IO\_cam\_sda} & K3 / J3 & 摄像头 I2C 配置总线 \\
\code{O\_cam\_24m} & L3 & 摄像头 24MHz 主时钟 \\
\code{O\_cam\_rst} & L2 & 摄像头复位 \\
\code{O\_tmds\_ch0/1/2\_p}, \code{O\_tmds\_clk\_p} & G16/F17/H14/C17 & HDMI TMDS \\
\code{I\_button[1:0]} & N2 / J14 & 板上按键（例程功能） \\
\code{O\_screen\_pwm} & R11 & 屏幕背光 PWM \\
\bottomrule
\end{tabular}
\end{center}

摄像头的数据 lane 走 MIPI D-PHY 硬核专用引脚（不在 pin.adc 中出现，
由 PHY IP 固定）。\textbf{这些约束已在工程里配好，不需要改。}

% ============================================================
\section{工程文件地图}

\subsection{目录结构}

\begin{Verilog}
08_board_bringup/01_video_passthrough/
├── lab_ex1/                          # 完整工程（整体拷贝即可用）
│   ├── td_project/
│   │   ├── camera_to_dsi_display.al  # TD 工程主文件(相对路径, 可搬移)
│   │   ├── camera_to_dsi_display.sdc # 时序约束
│   │   ├── al_ip/                    # 官方 IP
│   │   └── *_Runs/                   # 厂商历史运行报告(参考)
│   └── user_source/
│       ├── hdl_source/               # 96 个 RTL
│       │   ├── design_top_wrapper.v  # 顶层
│       │   ├── isp/ demosaic 等      # ISP 明码链
│       │   ├── mipi_dphy_rx/ hdmi_*  # PHY 与 HDMI(加密 .enc)
│       │   ├── ph1p35_ddr/           # DDR2 控制(官方)
│       │   └── uics520_cfg/          # SC520CS 配置(寄存器表+I2C 引擎)
│       └── constraints_source/       # pin.adc + timing.sdc
├── sim/                              # ModelSim 仿真(第 5 章)
└── README.md
\end{Verilog}

\subsection{SC520CS 适配要点（厂商《代码改动说明》+ 我们核对的）}

\begin{itemize}
  \item Lab1/2/3 已统一为 SC520CS \textbf{1280x720@60fps} 原厂 HD 模式；
  \item 摄像头配置引擎 \code{uicfgcs520} 上电先读传感器 ID，
    \textbf{0xEE4B 或 0x9E4B} 才算识别通过，然后写入
    \code{uics520reg\_720p60} 寄存器表（24MHz MCLK、2-lane、RAW10）；
  \item Lab1 另存有 \code{uics520reg\_full\_2560x1920\_30.v} 备用表
    （2560x1920@30，\textbf{未接入}工程，切换需同步改 MIPI/ISP/DDR/
    HDMI/时钟全套——官方 ROM 文件头有连锁清单）；
  \item \code{.al} 工程内全部是\textbf{相对路径}，整个 \code{lab\_ex1}
    文件夹拷到任何位置都能打开。
\end{itemize}

\begin{keybox}{分辨率变更的连锁清单（官方原文精神）}
换分辨率/帧率\textbf{不是换一张寄存器表}：MIPI 接收参数、图像宽高、
ISP、DDR 帧缓存、HDMI 时序、HDMI TX 参数、时钟与约束必须同组。
本项目锁定 720p60 不动；要改走接口变更流程（嵌入式教程（二））。
\end{keybox}

% ============================================================
\section{TD 导入与编译}

\subsection{环境}

TD 6.2.1（\code{D:\textbackslash TD6.2.1}，与工程创建版本 6.2.168116 相同）。
加密核（MIPI PHY / HDMI TX）随 TD 的 license 解密编译，无需额外操作。

\subsection{GUI 步骤（推荐首次使用）}

\begin{enumerate}
  \item 打开 TD $\rightarrow$ \code{File} $\rightarrow$ \code{Open Project}，
        选择工程文件（仓库 \code{08\_board\_bringup/} 下）：\\
        \code{01\_video\_passthrough/lab\_ex1/td\_project/}\\
        \code{camera\_to\_dsi\_display.al}；
  \item 打开后核对：器件 \code{PH1P35MDG324 -3}、顶层
        \code{design\_top\_wrapper}、源文件列表非空（约 57 项）、
        约束已挂（pin.adc + sdc）。若提示路径问题，用 Open Project
        重选 .al，TD 以 .al 所在目录为根重新解析相对路径；
  \item 双击流程窗口依次执行：\code{Compile}（综合）$\rightarrow$
        \code{Place \& Route}（布局布线，含时序分析）$\rightarrow$
        \code{Generate Bitstream}（出 bit）；
  \item 查看时序报告：\code{Report} $\rightarrow$ Timing Summary，
        \textbf{WNS $\geq$ 0} 为收敛。厂商历史 Runs 目录里有参考报告；
  \item 产物：\code{td\_project} 下 Runs 目录内的 \code{.bit}。
\end{enumerate}

\subsection{TCL 方式（可选，二次以后用）}

\begin{Verilog}
# TD 的 Tcl 控制台（命令名以本机 TD 6.2.1 补全为准）
prj_open camera_to_dsi_display.al
compile -impl syn_1
pnr     -impl syn_1
bitgen  -impl syn_1
rep_timing -impl syn_1
\end{Verilog}

\begin{warnbox}{三条导入纪律}
\textbf{1. 不要把 sim/ 下的任何文件加进 TD 工程}（替身会与真网表/
真原语抢模块名）；\textbf{2. 不要手改加密文件}（.enc 由 TD 解密）；
\textbf{3. 编译报"找不到文件"时先检查整个 lab\_ex1 文件夹是否完整拷贝}
（.al 用相对路径引用 \code{../user\_source/}，父目录缺文件即报错）。
\end{warnbox}

% ============================================================
\section{ModelSim 整链仿真（不需要板子）}

\subsection{目的与边界}

仿真回答一个问题：\textbf{"摄像头数据经过 ISP 后，画面长什么样、
能不能完整到达显示端"}。它覆盖的是\textbf{明码段}：
RAW10 输入 $\rightarrow$ demosaic $\rightarrow$ 128/96 重排 $\rightarrow$ AWB；
MIPI PHY、DDR2、HDMI 是加密/物理黑盒，不在仿真内
（HDMI 以"逐帧转储图片"替身验证）。

\subsection{一键运行}

\begin{Verilog}
cd 08_board_bringup\01_video_passthrough\sim
run_sim.bat                        :: 或在 ModelSim 控制台: do msim_passthrough.do
\end{Verilog}

期望输出（本仓库实测）：

\begin{Verilog}
# TB_RESULT: PASS frames_dumped=2 tuser_total=2 tlast_total=1370 ser_ovf=0
\end{Verilog}

出图（把仿真"画面"转成 PNG 肉眼看）：

\begin{Verilog}
python tools/dump2png.py 08_board_bringup/01_video_passthrough/sim passthrough_frame 0 1
\end{Verilog}

\subsection{仿真链路构成（为什么这样搭）}

\begin{Verilog}
raw10_beats.hex (虚拟传感器: BGGR 灰底+红块, 4 帧)
   |  1/8 占空馈送(真实 MIPI 节拍; 满速会淹没下游消费端)
   v
isp_top (官方明码副本: demosaic -> 128->96 -> AWB)
   |  96bit 抽头 (AWB 输出 = 视觉算法将来挂载的同一节点)
   v
px_serializer_fifo (仿真专用弹性串化器: 8192 beat FIFO + 真反压)
   |  1 像素/拍
   v
frame_dump24 -> passthrough_frame_N.hex -> dump2png -> PNG
\end{Verilog}

\begin{warnbox}{为什么不用 ser96to24（工程里的窃流串化器）}
它是\textbf{无反压}设计（13 深小 FIFO），吸不住 demosaic 的行级突发，
整链仿真里会丢 30\% 以上像素（帧永远不完整）。真实链路里这一级
由 DDR 帧缓存承担缓冲——仿真里用 \code{px\_serializer\_fifo}
充当"浅 DDR"，带真反压，实测 \code{ser\_ovf=0} 零丢拍。
\end{warnbox}

\subsection{结果判读}

\begin{center}
\begin{tabular}{lll}
\toprule
\rowcolor{rowtint}\textbf{字段} & \textbf{实测} & \textbf{含义} \\
\midrule
frames\_dumped & 2 & 完整转储帧数（每帧恰 921600 像素） \\
tuser\_total & 2 & 帧首标志数（每帧 1 次） \\
tlast\_total & 1370 & 行尾标志数（厂商 zhenghe 行窗口有合并， \\
& & 以\textbf{像素完整性}为准，不影响图像正确性） \\
ser\_ovf & 0 & 弹性串化器零拒收 \\
\bottomrule
\end{tabular}
\end{center}

注：厂商 demosaic 行窗对输出有节流（实测约 155ms/帧仿真时间），
TB 看门狗设 700ms 保证 2 帧完整；仿真跑完约 10 分钟属正常。

PNG 里应看到：灰色底（128 中灰）+ 0.2\% 确定性噪声点 +
中上部一个\textbf{红色条纹方块}（虚拟传感器的目标物；因厂商
demosaic 的 CFA 行相位约定，方块呈红/绿隔行——红色行已足够
检测链识别）。这就是"画面完整到达"的仿真证据；
真机色彩以实际 SC520CS 为准（其 CFA 相位由硬件决定，
天然匹配厂商链路）。

\begin{warnbox}{搭建这个仿真时踩掉的三个坑（复用价值最高）}
\textbf{1. 虚拟传感器向量为负值}：生成器用 numpy int32 做移位打包，
\code{p0<<30} 溢出——背景拍高 10bit 截断、噪声拍变负数，
\code{\$readmemh} 遇 \code{-} 非法字符\textbf{终止加载、其后全 X}。
修复：打包前先转 Python int（见 \code{gen\_raw10\_sensor.py} 注释）。\\
\textbf{2. ERAM 替身参数名不匹配}：官方 wrapper 传
\code{.ADDR\_WIDTH\_A(11)}，替身端口宽度却声明为
\code{ADDRA\_WIDTH(14)}——传入宽度没作用到端口，formal 高位悬空 Z，
\code{mem[Z 索引]} 读出 X 污染整帧。修复：参数名逐一对齐官方实例化
（替身 v3，见 \code{ph1p\_logic\_eram\_stub.v} 头注释）。\\
\textbf{3. TB 馈送违反 AXI 握手}：valid 在两个馈送槽之间保持高，
\textbf{同一 beat 被 ISP 重复消费 8 次}，eol/tuser 计数 x8 打乱
Bayer 行相位，颜色被抹平。修复：valid 置起 $\rightarrow$ ready 消费
$\rightarrow$ 立刻拉低，每拍恰好一次。\\[3pt]
三条都在嵌入式方向迭代记录（第 9、10 轮）有完整复盘。
\end{warnbox}

\begin{tipbox}{给后续算法联调的礼物}
96bit 抽头（AWB 输出）正是视觉检测链的挂载点。要预演算法效果，
把 detect\_top 挂进这个仿真即可离线调阈值——嵌入式教程（一）的
\code{tb\_system\_integration} 就是这么做的。
\end{tipbox}

% ============================================================
\section{烧录与上电验收}

\subsection{两种下载方式}

\begin{center}
\begin{tabular}{lll}
\toprule
\rowcolor{rowtint}\textbf{方式} & \textbf{工具} & \textbf{用途} \\
\midrule
JTAG 下载（SRAM） & TD 下载器 & 调试：快，断电即失 \\
QSPI 固化 & BitWizard & 交付：上电即运行（赛事要求） \\
\bottomrule
\end{tabular}
\end{center}

首次调试建议先 JTAG 下载（快），行为确认后再固化。

\subsection{JTAG 下载（调试模式）}

\begin{enumerate}
  \item 连接下载器与板上 JTAG 口，板上电；
  \item TD $\rightarrow$ \code{Tools} $\rightarrow$ \code{Download}
        （下载器界面）$\rightarrow$ 识别到 FPGA；
  \item 加载 .bit $\rightarrow$ \code{Download}（SRAM 模式）；
  \item 显示器应在数秒内出图。
\end{enumerate}

\subsection{QSPI 固化（上电即运行）}

\begin{enumerate}
  \item TD $\rightarrow$ \code{Tools} $\rightarrow$ \code{BitWizard}；
  \item 目标器件选\textbf{板载 QSPI Flash}（不是 FPGA SRAM）；
  \item 加载 .bit，用 Boot Image Tool 转换/直接烧写 Flash；
  \item 烧写完成 $\rightarrow$ \textbf{断电} $\rightarrow$ 拔下载器
        $\rightarrow$ 重新上电；
  \item 验收：显示器 15 秒内出图；再拔掉电脑（只留显示器与板子）
        重新上电仍出图——\textbf{"无 PC 参与"达标}。
\end{enumerate}

\subsection{上电五步联调节奏（每步过了才下一步）}

\begin{enumerate}
  \item 例程原样跑通出图（本步目标）；
  \item 挂 detect\_top（旁路）：画面无变化 + ChipWatcher 见 t\_valid；
  \item 挂 OSD：屏幕出现 TRK0 文字与目标框；
  \item 接舵机（独立 5V 共地）：duty 随目标移动；
  \item 固化 + 冷启动验收。
\end{enumerate}

% ============================================================
\section{故障排查表}

\begin{center}
\begin{tabular}{p{3.6cm}p{4.6cm}p{4.6cm}}
\toprule
\rowcolor{rowtint}\textbf{症状} & \textbf{先查} & \textbf{处置} \\
\midrule
完全黑屏 & 摄像头 I2C 配置是否完成 & 逻辑分析仪/ILA 抓 SCL/SDA；
读 ID 应为 EE4B/9E4B；查 FPC 方向与 O\_cam\_rst/24m \\
有数据无帧同步 & demosaic tuser 链 & 嵌入式教程（一）第 5 节的
界碑问题；确认 TD 仿真库环境与真机时序 \\
图像花屏/错位 & DDR 帧缓存 & 分辨率连锁清单是否被改动；DDR 引脚与
供电；复位时序 \\
颜色怪（偏色） & AWB 自适应 & 正常现象；演示前锁曝光或用阈值组切换 \\
HDMI 无信号 & 显示器与 TMDS & 显示器支持 720p60；换线；
查 VTC 时序参数 \\
时序不收敛（WNS<0） & 综合后报告 & 官方已知现象；
视觉逻辑挂载后重跑报告确认恶化量；关注 gimbal\_ctrl 帧拍单周期路径 \\
图像整体错行一行 & 综合报告 & 确认 bin\_smooth3x3 行缓存与 ser96to24
缓冲（均组合读）映射为 LUTRAM（误映射 BRAM 会错行；外脑评审转检查单） \\
\bottomrule
\end{tabular}
\end{center}

\begin{keybox}{黑屏排查的最短路径}
按信号流顺序二分：I2C 配置完成（ID 对）$\rightarrow$ 摄像头出流
（MIPI 波形）$\rightarrow$ ISP 输出（ChipWatcher 抓 O\_tvalid）
$\rightarrow$ DDR 读写 $\rightarrow$ HDMI。每级都有可抓的观测点，
从第一处断点定位问题——不要跳着猜。
\end{keybox}

% ============================================================
\section{下一步衔接}

\begin{itemize}
  \item 通路点亮后，按 \code{filelist.f}（嵌入式教程（二））把视觉
    13 个 RTL 加入工程，挂 detect\_top / osd\_track / vision\_top；
  \item 云台舵机的引脚与整定见嵌入式教程（二）与 Step3 教程；
  \item 所有接口位图以 \code{06\_interfaces/current/} 为唯一权威；
  \item 遇到本 PDF 未覆盖的问题，先查 \code{08\_board\_bringup/README.md}
    的"已知注意点"。
\end{itemize}

\begin{okbox}{验收清单（本步）}
\begin{itemize}
  \item [ ] 仿真：TB\_RESULT PASS，PNG 出图（红块清晰）
  \item [ ] TD 编译零错误，WNS $\geq$ 0
  \item [ ] JTAG 下载后显示器出图
  \item [ ] QSPI 固化后冷启动 15 秒内出图
  \item [ ] 拔 PC 重新上电仍出图
\end{itemize}
\end{okbox}

\end{document}
